\documentclass[a4paper,12pt]{article}
\usepackage[T2A]{fontenc}
\usepackage[utf8]{inputenc}
\usepackage[russian]{babel}
\usepackage{amsmath,amssymb,mathtools}
\usepackage[sfdefault,type1]{noto}
\usepackage{icomma}
\usepackage[a4paper,left=21mm,right=21mm,top=20mm,bottom=20mm]{geometry}
\usepackage{microtype}
\usepackage{tikz}
\usetikzlibrary{calc,arrows.meta,positioning,shapes.geometric,decorations.pathmorphing}
\usepackage{booktabs,tabularx}
\usepackage{enumitem}
\usepackage{fancyhdr}
\usepackage{setspace}
\usepackage{xcolor}
\usepackage{hyperref}
\hypersetup{hidelinks,pdftitle={Чек-лист: соединение пружин, трение и второй закон Ньютона},pdfauthor={Лёвин Артём Александрович}}
\definecolor{accent}{HTML}{16767A}
\definecolor{darkaccent}{HTML}{225256}
\definecolor{textgray}{HTML}{33393D}
\definecolor{muted}{HTML}{697276}
\color{textgray}
\setstretch{1.16}
\setlength{\parindent}{0pt}
\setlength{\parskip}{4pt}
\setlist[itemize]{leftmargin=7mm,itemsep=2pt,topsep=4pt,parsep=0pt}
\setlist[enumerate]{leftmargin=9mm,itemsep=9pt,topsep=6pt,parsep=0pt}
\newcommand{\key}[1]{\textcolor{darkaccent}{\textbf{#1}}}
\newcommand{\sect}[1]{\vspace{7pt}{\Large\bfseries\color{darkaccent}#1}\par\vspace{3pt}{\color{accent!65}\rule{\linewidth}{0.5pt}}\par\vspace{4pt}}
\newcommand{\subsect}[1]{\vspace{6pt}{\large\bfseries\color{darkaccent}#1}\par\vspace{2pt}}
\pagestyle{fancy}
\fancyhf{}
\fancyhead[L]{\small\color{darkaccent} Физика \textbullet\ Динамика}
\fancyhead[R]{\small\color{darkaccent} Володя Хачатурян}
\fancyfoot[C]{\small\color{muted}\thepage}
\renewcommand{\headrulewidth}{0.3pt}
\setlength{\headheight}{15pt}
\begin{document}
\hypersetup{pageanchor=false}
\begin{titlepage}
\thispagestyle{empty}
\begin{center}
\vspace*{24mm}
{\large\color{darkaccent} ФИЗИКА}\par
\vspace{13mm}
{\fontsize{28}{35}\selectfont\bfseries Чек-лист}\par
\vspace{12mm}
{\fontsize{21}{28}\selectfont\color{accent} Соединение пружин,\\[3mm] трение и второй закон Ньютона}\par
\vspace{11mm}
{\normalsize 9 класс \textbullet\ Закон Гука \textbullet\ Силы трения \textbullet\ Равнодействующая}\par
\vspace{26mm}
{\Large Володя Хачатурян}\par
\vspace{11mm}
{\large 31 октября 2026 г.}\par
\vfill
{\normalsize Лёвин Артём Александрович}\\[2mm]
{\small эксклюзивно для Володи Хачатуряна}\par
\vspace{21mm}
\end{center}
\begin{tikzpicture}[remember picture,overlay]
\draw[accent,line width=0.85pt]
  ([xshift=35mm,yshift=23mm]current page.south west)
  .. controls ([xshift=68mm,yshift=37mm]current page.south west) and
  ([xshift=90mm,yshift=12mm]current page.south west)
  .. ([xshift=114mm,yshift=23mm]current page.south west)
  .. controls ([xshift=141mm,yshift=36mm]current page.south west) and
  ([xshift=163mm,yshift=12mm]current page.south west)
  .. ([xshift=-35mm,yshift=23mm]current page.south east);
\end{tikzpicture}
\end{titlepage}
\hypersetup{pageanchor=true}

\sect{1. Как складываются жёсткости пружин}
Для одной пружины в области малых упругих деформаций закон Гука связывает \emph{модули} силы и изменения длины:
\[
\boxed{F_{\text{упр}}=k\,\Delta x},\qquad
\boxed{k=\frac{F_{\text{упр}}}{\Delta x}}.
\]
Жёсткость $k$ измеряется в \textbf{Н/м}; она является \textbf{скалярной величиной} и направления не имеет. Сила упругости направлена так, чтобы устранить деформацию.

\begin{center}
\begin{tikzpicture}[x=1cm,y=1cm,>=Latex]
\begin{scope}[shift={(0,0)}]
\draw[darkaccent,thick] (-.05,1.72)--(3.45,1.72);
\draw[darkaccent,thick] (-.05,-1.32)--(3.45,-1.32);
\foreach \x in {.75,2.55}{
  \draw[darkaccent,thick] (\x,1.72)--(\x,1.49);
  \draw[accent,thick,decorate,decoration={coil,aspect=.5,segment length=5pt,amplitude=2.7pt}] (\x,1.49)--(\x,-1.07);
  \draw[darkaccent,thick] (\x,-1.07)--(\x,-1.32);
}
\node[above,font=\small] at (.75,1.76) {$k_1$};\node[above,font=\small] at (2.55,1.76) {$k_2$};
\node[below,font=\small,align=center] at (1.7,-1.45) {Параллельно\\[1mm]$\Delta x_1=\Delta x_2$};
\end{scope}
\begin{scope}[shift={(7,0)}]
\draw[darkaccent,thick] (.1,1.72)--(3.35,1.72);
\draw[darkaccent,thick] (1.7,1.72)--(1.7,1.46);
\draw[accent,thick,decorate,decoration={coil,aspect=.5,segment length=5pt,amplitude=2.7pt}] (1.7,1.46)--(1.7,.18);
\draw[darkaccent,thick] (1.7,.18)--(1.7,-.05);
\draw[accent,thick,decorate,decoration={coil,aspect=.5,segment length=5pt,amplitude=2.7pt}] (1.7,-.05)--(1.7,-1.07);
\draw[darkaccent,thick] (1.7,-1.07)--(1.7,-1.32);
\node[right,font=\small] at (1.95,.81) {$k_1$};\node[right,font=\small] at (1.95,-.62) {$k_2$};
\node[below,font=\small,align=center] at (1.7,-1.45) {Последовательно\\[1mm]$F_1=F_2$};
\end{scope}
\end{tikzpicture}
\end{center}

\subsect{Параллельное соединение}
У пружин одинаковое удлинение. Общая сила равна сумме сил, действующих со стороны каждой пружины:
\[
F=F_1+F_2=(k_1+k_2)\Delta x
\quad\Longrightarrow\quad
\boxed{k_{\text{общ}}=k_1+k_2}.
\]
\key{Итог:} эквивалентная жёсткость больше каждой отдельной. При одинаковом удлинении систему растянуть труднее.

\subsect{Последовательное соединение}
Через каждую пружину передаётся одна и та же сила, а удлинения складываются:
\[
\Delta x=\frac{F}{k_1}+\frac{F}{k_2}
\quad\Longrightarrow\quad
\boxed{\frac{1}{k_{\text{общ}}}=\frac{1}{k_1}+\frac{1}{k_2}}.
\]
Для двух пружин:\quad$\displaystyle\boxed{k_{\text{общ}}=\frac{k_1k_2}{k_1+k_2}}$.
\key{Итог:} эквивалентная жёсткость меньше каждой отдельной.

\clearpage
\sect{2. Сравнение двух соединений}
\key{Пример.} Даны две пружины с жёсткостями $k_1=200$~Н/м и $k_2=300$~Н/м. К системе приложили растягивающую силу $F=12$~Н. Сравним удлинения в пределах действия закона Гука.

\subsect{Параллельно}
\[
k_{\text{пар}}=200+300=500\ \text{Н/м},\qquad
\Delta x_{\text{пар}}=\frac{12}{500}=0{,}024\ \text{м}=2{,}4\ \text{см}.
\]
У каждой пружины удлинение $2{,}4$~см. Но силы \emph{разные}:
\[
F_1=200\cdot0{,}024=4{,}8\ \text{Н},\quad
F_2=300\cdot0{,}024=7{,}2\ \text{Н};\quad F_1+F_2=12\ \text{Н}.
\]

\subsect{Последовательно}
\[
k_{\text{посл}}=\frac{200\cdot300}{200+300}=120\ \text{Н/м},\qquad
\Delta x_{\text{посл}}=\frac{12}{120}=0{,}1\ \text{м}=10\ \text{см}.
\]
Удлинения отдельных пружин: $\Delta x_1=12/200=6$~см, $\Delta x_2=12/300=4$~см. В сумме $10$~см.

\key{Быстрая самопроверка:} $120<200<300<500$. Последовательная система мягче, параллельная жёстче.

\sect{3. Сила трения: покой и скольжение}
\key{Трение покоя} возникает, когда поверхности соприкасаются и существует попытка сдвинуть их относительно друг друга. Оно препятствует началу относительного скольжения.

\key{Трение скольжения} действует, когда соприкасающиеся поверхности уже скользят друг по другу. Оно направлено против их \emph{относительной} скорости.

\begin{center}
\renewcommand{\arraystretch}{1.33}
\begin{tabularx}{\linewidth}{@{}>{\raggedright\arraybackslash}p{39mm}X@{}}
\toprule
Ситуация & Модуль силы трения \\
\midrule
Тело остаётся в покое & $0\le F_{\text{тр.пок}}\le F_{\text{тр.пок,max}}=\mu_{\text{п}}N$ \\
Тело скользит & $F_{\text{тр.ск}}=\mu_{\text{ск}}N$ (школьная модель) \\
\bottomrule
\end{tabularx}
\end{center}
Здесь $N$ --- модуль силы реакции опоры; $\mu_{\text{п}}$ и $\mu_{\text{ск}}$ --- безразмерные коэффициенты трения покоя и скольжения.

\textbf{Важно:} при покое трение подстраивается под приложенную силу до предела. Часто этот предел выше трения скольжения. Их равенство допустимо лишь при явном упрощении условия.

\clearpage
\sect{4. Когда тело начинает двигаться}
Пусть брусок первоначально покоится на горизонтальном столе. Мы тянем его вправо силой $F$, направленной горизонтально; других горизонтальных сил нет. Тогда $N=mg$.

\begin{itemize}
\item Если $F\le\mu_{\text{п}}N$, брусок может оставаться в покое; тогда $F_{\text{тр.пок}}=F$.
\item Если $F>\mu_{\text{п}}N$, равновесие невозможно и брусок начинает скользить. При скольжении используем $F_{\text{тр.ск}}=\mu_{\text{ск}}N$.
\end{itemize}

\subsect{Пример. Одна и та же поверхность, две силы}
Масса бруска $m=4$~кг; $\mu_{\text{п}}=0{,}25$, $\mu_{\text{ск}}=0{,}20$; $g=10$~Н/кг.
\[
N=mg=40\ \text{Н},\qquad
F_{\text{тр.пок,max}}=0{,}25\cdot40=10\ \text{Н}.
\]
\key{Случай 1.} Тянем с силой $6$~Н. Поскольку $6<10$, брусок покоится, сила трения равна $6$~Н. Ошибка --- автоматически подставлять сюда $\mu_{\text{ск}}N=8$~Н.

\key{Случай 2.} Тянем с силой $14$~Н. Порог в $10$~Н превышен, начинается скольжение. При его наличии $F_{\text{тр.ск}}=0{,}20\cdot40=8$~Н. Равнодействующая по горизонтали равна $14-8=6$~Н, ускорение $a=6/4=1{,}5$~м/с$^2$.

\subsect{Как меняется трение при изменении коэффициента}
Из $F_{\text{тр.ск}}=\mu_{\text{ск}}N$ при неизменном $N$ следует: коэффициент увеличили в $2$ раза $\Rightarrow$ силу трения скольжения увеличили в $2$ раза; уменьшили коэффициент в $3$ раза $\Rightarrow$ сила трения скольжения уменьшилась в $3$ раза.

\key{Условия приёма:} сравниваем одно и то же значение $N$ и один режим трения. Для трения покоя ниже предела такое рассуждение неприменимо: фактическая сила трения зависит от приложенной силы.

\sect{5. Сопротивление среды}
При движении тела в воздухе или жидкости может возникать сила сопротивления, направленная против относительной скорости тела и среды. Её величина, вообще говоря, зависит от скорости и свойств среды. Формулу $\mu N$ для сопротивления воздуха применять нельзя.

\clearpage
\sect{6. Второй закон Ньютона: тело на столе}
Рассмотрим брусок, который скользит \emph{вправо} по горизонтальному столу; к нему приложена горизонтальная сила тяги $F_{\text{т}}$. Сопротивлением воздуха пренебрегаем.

\begin{center}
\begin{tikzpicture}[>=Latex,scale=1]
\draw[darkaccent!45,line width=1.1pt] (-3,-1.15)--(4.0,-1.15);
\foreach \x in {-2.7,-2.25,...,3.8}{\draw[darkaccent!25] (\x,-1.15)--++(-.17,-.14);}
\draw[fill=accent!8,draw=darkaccent,thick,rounded corners=1pt] (-.85,-1.15) rectangle (.85,.05);
\coordinate (C) at (0,-.55);
\fill[darkaccent] (C) circle (1.5pt);
\draw[->,darkaccent,thick] (C)--++(0,1.85) node[above] {$\vec N$};
\draw[->,darkaccent,thick] (C)--++(0,-1.45) node[below] {$m\vec g$};
\draw[->,accent,very thick] (C)--++(3.2,0) node[right] {$\vec F_{\text{т}}$};
\draw[->,darkaccent,thick] (C)--++(-2.5,0) node[left] {$\vec F_{\text{тр}}$};
\draw[->,muted] (2.15,.87)--++(1.22,0) node[right] {$\vec a,\ \vec v$};
\draw[->,muted] (3.65,-1.8)--++(.48,0) node[right] {$x$};
\draw[->,muted] (3.65,-1.8)--++(0,.45) node[above] {$y$};
\end{tikzpicture}
\end{center}

\subsect{Три уравнения, которые нужно записать}
Направим ось $x$ вправо, ось $y$ вверх. Если вертикального ускорения нет и вертикально действуют лишь $N$ и $mg$, то
\[
\boxed{F_{\text{т}}-F_{\text{тр}}=ma},\qquad
\boxed{N-mg=0},\qquad
\boxed{F_{\text{тр}}=\mu_{\text{ск}}N}.
\]
Отсюда $N=mg$, $F_{\text{тр}}=\mu_{\text{ск}}mg$ и
\[
\boxed{F_{\text{т}}-\mu_{\text{ск}}mg=ma},\qquad
\boxed{F_{\text{т}}=m(a+\mu_{\text{ск}}g)}.
\]
\key{Знак минус} перед трением возникает потому, что эта сила направлена против положительной оси $x$.

\subsect{Пример. Найдём ускорение и равнодействующую}
Брусок массой $3$~кг скользит по столу вправо. На него горизонтально вправо действует сила тяги $12$~Н; $\mu_{\text{ск}}=0{,}2$, $g=10$~Н/кг.
\[
N=3\cdot10=30\ \text{Н},\quad F_{\text{тр}}=0{,}2\cdot30=6\ \text{Н};
\]
\[
F_{\text{рез},x}=12-6=6\ \text{Н},\qquad
\boxed{a=\frac{6}{3}=2\ \text{м/с}^2}.
\]
Силы $\vec N$ и $m\vec g$ взаимно компенсируются. Равнодействующая \emph{всех четырёх сил} направлена вправо и по модулю равна $6$~Н.

\clearpage
\begin{center}
{\LARGE\bfseries\color{darkaccent} Алгоритм: силы на горизонтальной поверхности}\par
\vspace{5mm}
{\small Для горизонтальной тяги, без других вертикальных сил и без вертикального ускорения}
\end{center}
\vspace{8mm}
\begin{center}
\begin{tikzpicture}[
 >=Latex,node distance=8mm,
 start/.style={ellipse,draw=accent,line width=.8pt,align=center,text width=44mm,minimum height=10mm,font=\small},
 flowstep/.style={rectangle,rounded corners=2mm,draw=accent,line width=.8pt,align=center,text width=118mm,minimum height=17mm,inner sep=2mm,font=\small},
 decision/.style={diamond,aspect=3,draw=accent,line width=.9pt,align=center,text width=44mm,inner sep=1mm,font=\small},
 smallstep/.style={rectangle,rounded corners=2mm,draw=accent,line width=.8pt,align=center,text width=48mm,minimum height=18mm,inner sep=2mm,font=\small},
 arr/.style={-{Latex[length=2.5mm]},draw=darkaccent,thick}]
\node[start] (a) {Начало};
\node[flowstep,below=of a] (b) {Нарисуйте брусок, $\vec F_{\text{т}}$, $\vec F_{\text{тр}}$, $\vec N$, $m\vec g$; выберите оси $x$ и $y$.};
\node[flowstep,below=of b] (c) {По вертикали: $N-mg=0$, значит $N=mg$.};
\node[decision,below=13mm of c] (d) {Есть скольжение?};
\node[smallstep,below left=16mm and 5mm of d] (e) {Если пока покоится: проверьте $F_{\text{т}}\le\mu_{\text{п}}N$.\ В покое $F_{\text{тр}}=F_{\text{т}}$.};
\node[smallstep,below right=16mm and 5mm of d] (f) {Тело скользит: \\[1mm] $F_{\text{тр}}=\mu_{\text{ск}}N$.};
\node[flowstep,below=41mm of d] (g) {По оси $x$: $F_{\text{т}}-F_{\text{тр}}=ma$. Найдите искомую величину и проверьте единицы.};
\node[start,below=of g] (h) {Ответ};
\draw[arr] (a)--(b);\draw[arr] (b)--(c);\draw[arr] (c)--(d);
\draw[arr] (d.west)-| node[above,pos=.33,font=\small]{нет} (e.north);
\draw[arr] (d.east)-| node[above,pos=.33,font=\small]{да} (f.north);
\draw[arr] (e.south)|-([xshift=-18mm,yshift=5mm]g.north)--([xshift=-18mm]g.north);
\draw[arr] (f.south)|-([xshift=18mm,yshift=5mm]g.north)--([xshift=18mm]g.north);
\draw[arr] (g)--(h);
\end{tikzpicture}
\end{center}
\vspace{5mm}
\textbf{Уточнение.} Если приложенная сила превосходит максимальное трение покоя, тело срывается с места; далее для режима скольжения используется правая ветвь. Если брусок покоится, ускорение равно нулю. На наклонной плоскости и при наклонной тяге уравнение для $N$ меняется.

\clearpage
\sect{7. Типичные ошибки и итоговая памятка}
\begin{itemize}
\item \key{Путают жёсткость и силу:} $k$ измеряется в Н/м и не имеет направления; $\vec F_{\text{упр}}$ --- вектор.
\item \key{Складывают жёсткости последовательно:} при последовательном соединении складываются \emph{удлинения}, при параллельном --- \emph{силы}.
\item \key{Путают модули сил в параллельных пружинах:} равны удлинения, но силы зависят от индивидуальных $k$.
\item \key{Любое трение приравнивают к $\mu N$:} для трения покоя эта формула даёт лишь \emph{максимум} при коэффициенте $\mu_{\text{п}}$.
\item \key{Считают $N=mg$ всегда:} равенство требует отсутствия ускорения по вертикали и других вертикальных составляющих сил.
\item \key{Забывают векторность:} на схеме отметьте направления всех четырёх сил и ускорения; по осям используйте знаки проекций.
\item \key{Игнорируют вертикальные силы в равнодействующей:} покажите явно, что $N-mg=0$; тогда результирующая действительно горизонтальна.
\end{itemize}

\subsect{Что определить в первую очередь}
\begin{center}
\renewcommand{\arraystretch}{1.38}
\begin{tabularx}{\linewidth}{@{}>{\raggedright\arraybackslash}p{53mm}X@{}}
\toprule
Признак задачи & Полезный первый шаг \\
\midrule
Пружины рядом, одинаковое $\Delta x$ & $k_{\text{общ}}=k_1+k_2$ \\
Пружины друг за другом, одна сила & $\frac1{k_{\text{общ}}}=\frac1{k_1}+\frac1{k_2}$ \\
Тело ещё покоится & Сравнить тягу с $\mu_{\text{п}}N$ \\
Тело скользит & Найти $F_{\text{тр}}=\mu_{\text{ск}}N$ \\
Дано или требуется ускорение & Расписать $\sum F_x=ma$, $\sum F_y=0$ \\
Нужно найти равнодействующую & Учесть \emph{все} силы, затем $\vec F_{\text{рез}}=m\vec a$ \\
\bottomrule
\end{tabularx}
\end{center}

\subsect{Единицы и проверка}
\[
1\ \text{см}=0{,}01\ \text{м},\qquad
[k]=\text{Н/м},\qquad
[F]=\text{Н},\qquad [a]=\text{м/с}^2.
\]
В задачах ниже принимайте $g=10$~Н/кг. Если не указано иное, пружины работают в области действия закона Гука, а поверхность горизонтальна.

\clearpage
\sect{8. Тренировочный блок. Путь А}
\textcolor{muted}{\small 10 заданий по нарастающей сложности. Для задач на движение нарисуйте схему сил, укажите оси и запишите уравнения по осям.}
\begin{enumerate}[label=\textbf{\arabic*.}]
\item Две пружины с жёсткостями $160$ и $240$~Н/м соединены \textbf{параллельно}. Найдите эквивалентную жёсткость.
\item Две пружины с жёсткостями $120$ и $180$~Н/м соединены \textbf{последовательно}. Найдите эквивалентную жёсткость.
\item Пружины жёсткостью $300$ и $600$~Н/м соединяют сначала параллельно, затем последовательно. В каждом случае к системе прикладывают растягивающую силу $18$~Н. Найдите общее удлинение в обоих случаях.
\item Две пружины жёсткостью $200$ и $500$~Н/м, соединённые параллельно, растянули на $3$~см. Найдите силу упругости каждой пружины и общую силу.
\item Брусок массой $5$~кг скользит по горизонтальному столу. Коэффициент трения скольжения равен $0{,}3$. Найдите реакцию опоры и силу трения.
\item К первоначально покоящемуся бруску массой $2$~кг прикладывают горизонтальную силу $6$~Н. Коэффициент трения покоя равен $0{,}4$. Останется ли брусок в покое? Чему равна фактическая сила трения?
\item К тому же первоначально покоящемуся бруску массой $2$~кг приложили горизонтальную силу $12$~Н. Коэффициенты трения покоя и скольжения равны соответственно $0{,}4$ и $0{,}25$. Определите, начнётся ли движение, и найдите ускорение при скольжении.
\end{enumerate}

\clearpage
\sect{Тренировочный блок. Продолжение}
\begin{enumerate}[label=\textbf{\arabic*.},start=8]
\item Брусок массой $3$~кг скользит вправо по горизонтальному столу. На него действует горизонтальная сила тяги $12$~Н вправо; коэффициент трения скольжения равен $0{,}2$. Найдите модуль равнодействующей всех сил, её направление и ускорение бруска.
\item Брусок массой $4$~кг скользит по столу с ускорением $1{,}5$~м/с$^2$ вправо. Коэффициент трения скольжения равен $0{,}25$. Найдите необходимую горизонтальную силу тяги.
\item \textbf{Обобщение.} Две одинаково растянутые параллельные пружины жёсткостью $200$ и $300$~Н/м тянут брусок массой $6$~кг вправо по горизонтальной поверхности. Удлинение каждой пружины равно $4$~см. Брусок уже скользит вправо; коэффициент трения скольжения равен $0{,}2$. Считая, что других горизонтальных сил нет, найдите силу тяги, силу трения и ускорение бруска.
\end{enumerate}

\vspace{12mm}
\begin{center}
\begin{tikzpicture}[>=Latex]
\draw[darkaccent!45,line width=.9pt] (-5,0)--(5,0);
\node[anchor=west,align=left,text=muted,font=\small] at (-5,-.55) {Проверьте после решения: единицы, знаки проекций,\\[1mm] режим трения и физический смысл результата.};
\end{tikzpicture}
\end{center}

\clearpage
\sect{9. Ответы к тренировочному блоку}
\vspace{6mm}
\renewcommand{\arraystretch}{1.75}
\begin{center}
\begin{tabularx}{\linewidth}{@{}>{\bfseries}p{9mm}X@{}}
\toprule
№ & Ответ \\
\midrule
1 & $400$~Н/м. \\
2 & $72$~Н/м. \\
3 & Параллельно: $2$~см; последовательно: $9$~см. \\
4 & $F_1=6$~Н, $F_2=15$~Н, $F_{\text{общ}}=21$~Н. \\
5 & $N=50$~Н; $F_{\text{тр.ск}}=15$~Н. \\
6 & Останется в покое; $F_{\text{тр.пок}}=6$~Н. \\
7 & Движение начнётся; $F_{\text{тр.ск}}=5$~Н; $a=3{,}5$~м/с$^2$ вправо. \\
8 & Равнодействующая $6$~Н вправо; $a=2$~м/с$^2$ вправо. \\
9 & $F_{\text{т}}=16$~Н. \\
10 & $F_{\text{т}}=20$~Н, $F_{\text{тр.ск}}=12$~Н, $a=\dfrac43$~м/с$^2$ вправо. \\
\bottomrule
\end{tabularx}
\end{center}
\vfill
\begin{center}
{\small\color{muted}Лёвин Артём Александрович --- эксклюзивно для Володи Хачатуряна}\par
\end{center}
\end{document}
